% ===== PRÉAMBULE CANONIQUE — à coller VERBATIM en tête de CHAQUE .tex =====
\documentclass[11pt,a4paper]{article}
\usepackage[utf8]{inputenc}\usepackage[T1]{fontenc}\usepackage{lmodern}
\usepackage[french]{babel}
\usepackage{amsmath,amssymb,mathtools}\usepackage{icomma}
\usepackage{siunitx}\sisetup{locale=FR}  % \qty{}{}, \si{}, \ang{}
\usepackage{xcolor}\usepackage{colortbl}
\usepackage{tikz}\usepackage{pgfplots}\pgfplotsset{compat=1.18}
\usepackage[siunitx,europeanresistors]{circuitikz} % circuits (élec.)
\usepackage[most]{tcolorbox}\usepackage{enumitem}\usepackage{array,booktabs}
\usepackage[a4paper,margin=2.2cm]{geometry}
\usetikzlibrary{arrows.meta,calc,angles,quotes,positioning,patterns,%
  backgrounds,shapes.geometric,decorations.pathmorphing,decorations.markings,intersections}
\definecolor{primary}{RGB}{222,49,99}\definecolor{primarydark}{RGB}{150,22,55}
\definecolor{defcol}{RGB}{0,114,178}\definecolor{methcol}{RGB}{52,92,148}
\definecolor{excol}{RGB}{0,135,90}\definecolor{attncol}{RGB}{200,40,55}
\tcbset{boxbase/.style={enhanced,breakable,boxrule=0.9pt,arc=2.5pt,
  left=3mm,right=3mm,top=2.3mm,bottom=2.3mm,fonttitle=\bfseries,coltitle=white}}
% --- boîtes TUTORIEL (noms EXACTS, ne pas renommer) ---
\newtcolorbox{cle}[1][]{boxbase,colback=defcol!6,colframe=defcol,
  title={\textsc{À retenir}\ifx\relax#1\relax\relax\else\textmd{ : \itshape #1}\fi}}
\newtcolorbox{methode}[1][]{boxbase,colback=methcol!7,colframe=methcol,sharp corners,
  title={\textsc{Méthode}\ifx\relax#1\relax\relax\else\textmd{ --- \itshape #1}\fi}}
\newtcolorbox{exemplebox}[1][]{boxbase,colback=excol!5,colframe=excol,
  title={\textsc{Exemple}\ifx\relax#1\relax\relax\else\textmd{ : \itshape #1}\fi}}
\newtcolorbox{piege}[1][]{boxbase,colback=attncol!6,colframe=attncol,
  title={\textsc{Piège}\ifx\relax#1\relax\relax\else\textmd{ : \itshape #1}\fi}}
\newtcolorbox{remarque}[1][]{boxbase,colback=black!4,colframe=black!45,
  title={\textsc{Remarque}\ifx\relax#1\relax\relax\else\textmd{ : \itshape #1}\fi}}
% --- boîtes EXERCICE / CORRIGÉ (noms EXACTS, auto-comptées) ---
\newtcolorbox[auto counter]{exo}[1][]{boxbase,colback=primary!4,colframe=primary,
  title={\textsc{Exercice}~\thetcbcounter\ifx\relax#1\relax\relax\else\textmd{ --- \itshape #1}\fi}}
\newtcolorbox[auto counter]{corrige}[1][]{boxbase,colback=excol!4,colframe=excol,
  title={\textsc{Corrigé}~\thetcbcounter\ifx\relax#1\relax\relax\else\textmd{ --- \itshape #1}\fi}}
\newcommand{\vv}[1]{\vec{#1}}\newcommand{\ds}{\displaystyle}
\newcommand{\R}{\mathbb{R}}\newcommand{\N}{\mathbb{N}}\newcommand{\Z}{\mathbb{Z}}
% (un \usepackage{fancyhdr}+en-tête est possible ; il est ignoré côté web)
% =========================================================================
\begin{document}
% THEME_TITLE: Tir oblique (mouvement balistique)
\sisetup{per-mode=symbol}

\noindent Un projectile lancé en oblique combine deux mouvements indépendants : un \textbf{MRU}
horizontal et une \textbf{chute libre} verticale. Tout se déduit de la décomposition de $\vv{v}_0$.

\begin{cle}[décomposition de la vitesse initiale]
Avec un angle de tir $\theta$ par rapport à l'horizontale :
\[ v_{0x}=v_0\cos\theta\ \ (\text{MRU}),\qquad v_{0y}=v_0\sin\theta\ \ (\text{soumise à }g). \]
Selon $y$ : montée décélérée ($a=-g$, MRUD) puis descente accélérée ($a=+g$, MRUA).
\begin{center}
\begin{tikzpicture}[scale=0.85,>=Stealth]
  \draw[->,thick,black!70] (-0.2,0)--(8.2,0) node[right]{$x$};
  \draw[->,thick,black!70] (0,-0.2)--(0,3.2) node[above]{$y$};
  \draw[primary,line width=1.1pt,domain=0:7,smooth,variable=\x] plot ({\x},{0.22*\x*(7-\x)});
  \draw[->,line width=1pt,primary] (0,0)--(1.3,1.55) node[above right]{$\vv{v}_0$};
  \draw[->,line width=0.9pt,defcol] (0,0)--(1.3,0) node[midway,below]{$v_{0x}$};
  \draw[->,line width=0.9pt,attncol] (0,0)--(0,1.55) node[midway,left]{$v_{0y}$};
  \draw[primary,thick] (0.6,0) arc (0:50:0.6); \node[primary,font=\scriptsize] at (0.85,0.24){$\theta$};
  \fill[attncol] (3.5,2.7) circle (2pt);
  \draw[->,line width=0.9pt,defcol] (3.5,2.7)--(4.4,2.7); \node[attncol,font=\scriptsize,above] at (3.5,2.85){$v_y=0$};
  \draw[line width=1.5pt,excol] (0,-0.35)--(7,-0.35); \node[excol,font=\scriptsize] at (3.5,-0.62){portée $x_p$};
  \draw[dashed,black!35] (3.5,0)--(3.5,2.7);
\end{tikzpicture}
\end{center}
\end{cle}

\begin{cle}[sommet et hauteur maximale]
Au sommet la vitesse verticale s'annule ($v_y=0$) :
\[ t_s=\frac{v_{0y}}{g}=\frac{v_0\sin\theta}{g},\qquad
   y_{\max}=\frac{v_{0y}^2}{2g}=\frac{v_0^2\sin^2\theta}{2g}. \]
Attention : au sommet la vitesse n'est pas nulle, il reste $v=v_{0x}=v_0\cos\theta$ (horizontale).
\end{cle}

\begin{cle}[durée, portée et angle optimal]
Par symétrie (départ et arrivée à la même hauteur), la durée totale est $T=2t_s$ et la portée :
\[ x_p=v_{0x}\cdot 2t_s=\frac{v_0^2\sin(2\theta)}{g}\qquad\text{(maximale pour }\theta=\ang{45}). \]
\end{cle}

\begin{methode}[résoudre un tir oblique]
\begin{enumerate}[label=\arabic*.,leftmargin=1.6em,topsep=1pt,itemsep=0pt]
  \item Décompose : $v_{0x}=v_0\cos\theta$, $v_{0y}=v_0\sin\theta$.
  \item Sommet : $t_s=v_{0y}/g$, puis $y_{\max}=v_{0y}^2/2g$.
  \item Durée totale $T=2t_s$, portée $x_p=v_{0x}\,T$.
\end{enumerate}
\end{methode}

\begin{exemplebox}[bille lancée à \ang{30}]
$v_0=\qty{20}{\metre\per\second}$, $\theta=\ang{30}$, $g=\qty{10}{\metre\per\second\squared}$ :
$v_{0x}=10\sqrt3$, $v_{0y}=\qty{10}{\metre\per\second}$. Alors $t_s=\qty{1}{\second}$,
$y_{\max}=\qty{5}{\metre}$, $x_p=v_{0x}\times2t_s=20\sqrt3\approx\qty{34,6}{\metre}$.
\end{exemplebox}

\begin{piege}[deux idées fausses à éviter]
\textbf{(1)} Au sommet, la vitesse \emph{n'est pas nulle} : seule sa composante verticale l'est.\quad
\textbf{(2)} Deux angles \emph{complémentaires} (par ex. \ang{30} et \ang{60}) donnent la \emph{même portée},
car $\sin(2\theta)$ est identique.
\end{piege}

\begin{remarque}[une seule accélération]
Sur toute la trajectoire, l'accélération est constante : $\vv{a}=\vv{g}$ (verticale, vers le bas).
La composante horizontale du mouvement reste un MRU : le déplacement horizontal croît \emph{linéairement} avec le temps.
\end{remarque}

\end{document}
